\noindent\textbf{源码生成说明：} 本附录由冻结提交源码清单自动生成，共列入 24 个程序；代码内容与正式提交源码逐字节一致。当前仅因组别与参赛编号未填写而标记为 INTERNAL-QA。\par
\vspace{0.4em}
\Needspace{6\baselineskip}
\subsection{附录 C.1 项目级复现入口}
\noindent{}原项目相对路径：\texttt{run\_pipeline.py}\par
\vspace{0.25em}
{\fontsize{6.4pt}{7.2pt}\selectfont
\VerbatimInput[breaklines=true,breakanywhere=true,numbers=left,numbersep=5pt,frame=single,framesep=2mm,rulecolor=\color{black!35}]{source_appendix/001.py}
}
\Needspace{6\baselineskip}
\subsection{附录 C.2 附件分段重建与审计程序}
\noindent{}原项目相对路径：\texttt{02\_数据与参数/src/reconstruct\_segments.py}\par
\vspace{0.25em}
{\fontsize{6.4pt}{7.2pt}\selectfont
\VerbatimInput[breaklines=true,breakanywhere=true,numbers=left,numbersep=5pt,frame=single,framesep=2mm,rulecolor=\color{black!35}]{source_appendix/002.py}
}
\Needspace{6\baselineskip}
\subsection{附录 C.3 问题一有限圆柱接触与导通判定程序}
\noindent{}原项目相对路径：\texttt{问题/问题1/src/solve.py}\par
\vspace{0.25em}
{\fontsize{6.4pt}{7.2pt}\selectfont
\VerbatimInput[breaklines=true,breakanywhere=true,numbers=left,numbersep=5pt,frame=single,framesep=2mm,rulecolor=\color{black!35}]{source_appendix/003.py}
}
\Needspace{6\baselineskip}
\subsection{附录 C.4 问题二固定样本蒙特卡洛程序}
\noindent{}原项目相对路径：\texttt{问题/问题2/src/solve.py}\par
\vspace{0.25em}
{\fontsize{6.4pt}{7.2pt}\selectfont
\VerbatimInput[breaklines=true,breakanywhere=true,numbers=left,numbersep=5pt,frame=single,framesep=2mm,rulecolor=\color{black!35}]{source_appendix/004.py}
}
\Needspace{6\baselineskip}
\subsection{附录 C.5 问题三两阶段概率确认程序}
\noindent{}原项目相对路径：\texttt{问题/问题3/src/solve.py}\par
\vspace{0.25em}
{\fontsize{6.4pt}{7.2pt}\selectfont
\VerbatimInput[breaklines=true,breakanywhere=true,numbers=left,numbersep=5pt,frame=single,framesep=2mm,rulecolor=\color{black!35}]{source_appendix/005.py}
}
\Needspace{6\baselineskip}
\subsection{附录 C.6 问题四混合介质成本优化程序}
\noindent{}原项目相对路径：\texttt{问题/问题4/src/solve.py}\par
\vspace{0.25em}
{\fontsize{6.4pt}{7.2pt}\selectfont
\VerbatimInput[breaklines=true,breakanywhere=true,numbers=left,numbersep=5pt,frame=single,framesep=2mm,rulecolor=\color{black!35}]{source_appendix/006.py}
}
\Needspace{6\baselineskip}
\subsection{附录 C.7 有限平底圆柱几何核}
\noindent{}原项目相对路径：\texttt{公共代码/geometry\_kernel.py}\par
\vspace{0.25em}
{\fontsize{6.4pt}{7.2pt}\selectfont
\VerbatimInput[breaklines=true,breakanywhere=true,numbers=left,numbersep=5pt,frame=single,framesep=2mm,rulecolor=\color{black!35}]{source_appendix/007.py}
}
\Needspace{6\baselineskip}
\subsection{附录 C.8 几何窄相加速核}
\noindent{}原项目相对路径：\texttt{公共代码/fast\_geometry.py}\par
\vspace{0.25em}
{\fontsize{6.4pt}{7.2pt}\selectfont
\VerbatimInput[breaklines=true,breakanywhere=true,numbers=left,numbersep=5pt,frame=single,framesep=2mm,rulecolor=\color{black!35}]{source_appendix/008.py}
}
\Needspace{6\baselineskip}
\subsection{附录 C.9 单介质随机微构体仿真核}
\noindent{}原项目相对路径：\texttt{公共代码/microstructure\_sim.py}\par
\vspace{0.25em}
{\fontsize{6.4pt}{7.2pt}\selectfont
\VerbatimInput[breaklines=true,breakanywhere=true,numbers=left,numbersep=5pt,frame=single,framesep=2mm,rulecolor=\color{black!35}]{source_appendix/009.py}
}
\Needspace{6\baselineskip}
\subsection{附录 C.10 混合介质随机微构体仿真核}
\noindent{}原项目相对路径：\texttt{公共代码/mixed\_microstructure\_sim.py}\par
\vspace{0.25em}
{\fontsize{6.4pt}{7.2pt}\selectfont
\VerbatimInput[breaklines=true,breakanywhere=true,numbers=left,numbersep=5pt,frame=single,framesep=2mm,rulecolor=\color{black!35}]{source_appendix/010.py}
}
\Needspace{6\baselineskip}
\subsection{附录 C.11 二维资源连通与 Pareto 前沿算法}
\noindent{}原项目相对路径：\texttt{公共代码/pareto\_connectivity.py}\par
\vspace{0.25em}
{\fontsize{6.4pt}{7.2pt}\selectfont
\VerbatimInput[breaklines=true,breakanywhere=true,numbers=left,numbersep=5pt,frame=single,framesep=2mm,rulecolor=\color{black!35}]{source_appendix/011.py}
}
\Needspace{6\baselineskip}
\subsection{附录 C.12 结果注册与 LaTeX 导出程序}
\noindent{}原项目相对路径：\texttt{公共代码/result\_registry.py}\par
\vspace{0.25em}
{\fontsize{6.4pt}{7.2pt}\selectfont
\VerbatimInput[breaklines=true,breakanywhere=true,numbers=left,numbersep=5pt,frame=single,framesep=2mm,rulecolor=\color{black!35}]{source_appendix/012.py}
}
\Needspace{6\baselineskip}
\subsection{附录 C.13 科学绘图样式模块}
\noindent{}原项目相对路径：\texttt{公共代码/plot\_style.py}\par
\vspace{0.25em}
{\fontsize{6.4pt}{7.2pt}\selectfont
\VerbatimInput[breaklines=true,breakanywhere=true,numbers=left,numbersep=5pt,frame=single,framesep=2mm,rulecolor=\color{black!35}]{source_appendix/013.py}
}
\Needspace{6\baselineskip}
\subsection{附录 C.14 问题一三维场景数据生成程序}
\noindent{}原项目相对路径：\texttt{论文/figures/src/build\_q1\_scenes.py}\par
\vspace{0.25em}
{\fontsize{6.4pt}{7.2pt}\selectfont
\VerbatimInput[breaklines=true,breakanywhere=true,numbers=left,numbersep=5pt,frame=single,framesep=2mm,rulecolor=\color{black!35}]{source_appendix/014.py}
}
\Needspace{6\baselineskip}
\subsection{附录 C.15 FreeCAD 参数化实体建模程序}
\noindent{}原项目相对路径：\texttt{论文/figures/src/build\_freecad\_scene.py}\par
\vspace{0.25em}
{\fontsize{6.4pt}{7.2pt}\selectfont
\VerbatimInput[breaklines=true,breakanywhere=true,numbers=left,numbersep=5pt,frame=single,framesep=2mm,rulecolor=\color{black!35}]{source_appendix/015.py}
}
\Needspace{6\baselineskip}
\subsection{附录 C.16 FreeCAD 通用视图渲染程序}
\noindent{}原项目相对路径：\texttt{论文/figures/src/render\_freecad\_scene.py}\par
\vspace{0.25em}
{\fontsize{6.4pt}{7.2pt}\selectfont
\VerbatimInput[breaklines=true,breakanywhere=true,numbers=left,numbersep=5pt,frame=single,framesep=2mm,rulecolor=\color{black!35}]{source_appendix/016.py}
}
\Needspace{6\baselineskip}
\subsection{附录 C.17 问题二与问题三概率图生成程序}
\noindent{}原项目相对路径：\texttt{论文/figures/src/build\_q2\_q3\_figure.py}\par
\vspace{0.25em}
{\fontsize{6.4pt}{7.2pt}\selectfont
\VerbatimInput[breaklines=true,breakanywhere=true,numbers=left,numbersep=5pt,frame=single,framesep=2mm,rulecolor=\color{black!35}]{source_appendix/017.py}
}
\Needspace{6\baselineskip}
\subsection{附录 C.18 问题四成本前沿图生成程序}
\noindent{}原项目相对路径：\texttt{论文/figures/src/build\_q4\_frontier.py}\par
\vspace{0.25em}
{\fontsize{6.4pt}{7.2pt}\selectfont
\VerbatimInput[breaklines=true,breakanywhere=true,numbers=left,numbersep=5pt,frame=single,framesep=2mm,rulecolor=\color{black!35}]{source_appendix/018.py}
}
\Needspace{6\baselineskip}
\subsection{附录 C.19 问题四真实随机试验三维建模程序}
\noindent{}原项目相对路径：\texttt{论文/figures/src/build\_q4\_mixed\_scene.py}\par
\vspace{0.25em}
{\fontsize{6.4pt}{7.2pt}\selectfont
\VerbatimInput[breaklines=true,breakanywhere=true,numbers=left,numbersep=5pt,frame=single,framesep=2mm,rulecolor=\color{black!35}]{source_appendix/019.py}
}
\Needspace{6\baselineskip}
\subsection{附录 C.20 问题四三维场景渲染与审计程序}
\noindent{}原项目相对路径：\texttt{论文/figures/src/render\_q4\_mixed\_scene.py}\par
\vspace{0.25em}
{\fontsize{6.4pt}{7.2pt}\selectfont
\VerbatimInput[breaklines=true,breakanywhere=true,numbers=left,numbersep=5pt,frame=single,framesep=2mm,rulecolor=\color{black!35}]{source_appendix/020.py}
}
\Needspace{6\baselineskip}
\subsection{附录 C.21 问题四导通见证与有序接触图生成程序}
\noindent{}原项目相对路径：\texttt{论文/figures/src/build\_q4\_witness\_figure.py}\par
\vspace{0.25em}
{\fontsize{6.4pt}{7.2pt}\selectfont
\VerbatimInput[breaklines=true,breakanywhere=true,numbers=left,numbersep=5pt,frame=single,framesep=2mm,rulecolor=\color{black!35}]{source_appendix/021.py}
}
\Needspace{6\baselineskip}
\subsection{附录 C.22 问题四正式三维资产串行生成与总审计程序}
\noindent{}原项目相对路径：\texttt{论文/figures/src/build\_q4\_final\_assets.py}\par
\vspace{0.25em}
{\fontsize{6.4pt}{7.2pt}\selectfont
\VerbatimInput[breaklines=true,breakanywhere=true,numbers=left,numbersep=5pt,frame=single,framesep=2mm,rulecolor=\color{black!35}]{source_appendix/022.py}
}
\Needspace{6\baselineskip}
\subsection{附录 C.23 问题四正式确认完整整数域独立审计程序}
\noindent{}原项目相对路径：\texttt{问题/问题4/src/audit\_confirmation.py}\par
\vspace{0.25em}
{\fontsize{6.4pt}{7.2pt}\selectfont
\VerbatimInput[breaklines=true,breakanywhere=true,numbers=left,numbersep=5pt,frame=single,framesep=2mm,rulecolor=\color{black!35}]{source_appendix/023.py}
}
\Needspace{6\baselineskip}
\subsection{附录 C.24 统一建模流程与验证诊断图生成程序}
\noindent{}原项目相对路径：\texttt{论文/figures/src/build\_explanatory\_figures.py}\par
\vspace{0.25em}
{\fontsize{6.4pt}{7.2pt}\selectfont
\VerbatimInput[breaklines=true,breakanywhere=true,numbers=left,numbersep=5pt,frame=single,framesep=2mm,rulecolor=\color{black!35}]{source_appendix/024.py}
}
